% Hand-written interpretation of the generated numbers (nce_macros.tex and nce_table_*.tex).
\newcommand{\nceSiAccuracySentence}{its offline error and closed-loop success (\nceSiChunk{} chunk MSE, \nceSiSucc{} success) stay within the variation we measure across training seeds.}

\newcommand{\nceNeuroText}{%
\textbf{The core is cheap; the camera is not.} Table~\ref{tab:energy} breaks one control tick into components. The camera ViT accounts for \nceVitSharePct\% of the full model's dense operations, the LiDAR/radar pillar encoder for most of the rest, and the liquid--spiking core ($T{=}8$ steps, two layers) for under 1\%. Any claim of neuromorphic efficiency for the \emph{system} therefore depends first on the perception front-end, not on the core. The radar voxel grid is more than 99\% empty in our scenes, which makes the pillar encoder a natural candidate for event-driven or sparse convolution.

\textbf{Where the spikes go matters more than whether a network spikes.} Table~\ref{tab:neuro} compares four cores. In the as-built core, the spike trains of Q, K and V turn $QK^\top$ and $AV$ into accumulates. At the measured firing rate of \nceMeanRate{} these products become almost free, but in a short control window they are only \nceAttnFracPct\% of the core's operations. The projections $W_{qkv}$, $W_x$, $W_h$ and the output layers dominate, so the as-built saving over a dense core is \nceSavingAsBuiltPct\%. Adding $\Delta t$-aware LIF populations in front of $W_{qkv}$ and $W_x$ (spike-driven projection inputs) makes \nceSiFracPct\% of the core's operations spike-driven and saves \nceSiSavingPct\% of its energy, at a similar firing rate (\nceSiRate). The remaining dense operations are the recurrent $W_h h_{t-1}$, the attention output projection and the LTC read-out, whose inputs are real-valued states. Spiking those as well bounds the saving at about 98\% for these firing rates, but it would change the continuous-time state that makes the LTC useful, and we leave it to future work.

\textbf{Accuracy of the variants.} Across three training seeds of the as-built model, held-out chunk error is \nceSeedMse{} and closed-loop success \nceSeedSucc{} (per seed: \nceSeedList). Every variant in Table~\ref{tab:neuro} lies within that spread on both measures. We therefore do not claim that the LIF attention or the LTC recurrence \emph{improves} accuracy at this scale. The feed-forward (no-LTC) core has the lowest success and the highest error, but by one episode, which is within noise. The claim we do make is narrower and robust to this spread: \emph{moving the spikes to the projection inputs cuts the core's estimated energy by more than half without a detectable loss in accuracy.}}

\newcommand{\nceClosedText}{%
Table~\ref{tab:closed} evaluates the controllers on identical seeds. Verify-behind matches the full model's closed-loop success and collision rate while its critical path (radar pillars, the drafter latent, and the drafter core on a cache miss) costs \nceSpecCritUJ\,$\mu$J per tick, against \nceTargetUJ\,$\mu$J for running the full model every tick (\nceCritRatio$\times$ less), and 4.1 against 24.2\,ms of latency. Total energy is \emph{not} reduced, because the target still verifies every window off the critical path; the gain is in what the control loop waits for. The drafter alone collides in 3\% of episodes even with the CBF filter, and adding verification removes those collisions. The target accepted 75\% of executed drafts. A muscle-memory cache of verified chunks supplied most actions, which shortens the critical path but raises action jerk as the source switches between cache, drafter and override. Latency was measured on a shared 2-core CPU with the action head evaluated over the whole window; the energy figures assume the head is decoded at the last step only, which gives identical actions.}

\newcommand{\nceDiagText}{%
Table~\ref{tab:diag} collects the failure modes from \S\ref{sec:training}. The copycat policy had the \emph{lowest} held-out imitation error of all runs and 0\% success; the colour-blind policy had a similar error and 4\% success. Only closed-loop evaluation revealed either. Two rounds of DAgger, in which the policy flies and the expert labels the states it actually visits, halved the mean final distance to the goal without raising success. This points to demonstration coverage and expert smoothness as the main limits on absolute performance, rather than the neuromorphic core.}

\newcommand{\nceDiscussionText}{%
Three messages carry over beyond this system. First, for spiking control policies, operation-level accounting should accompany any efficiency claim. A spiking-transformer core can be spiking in name while less than 2\% of its operations are spike-driven. Second, the perception front-end dominates a multimodal policy's energy, so neuromorphic gains at system level require event-driven perception (event cameras, sparse radar/LiDAR processing), not only a spiking core. Third, the continuous-time components integrate the measured $\Delta t$ at no extra operation cost: both the LIF leak and the closed-form LTC update are elementwise. They also remain stable under irregular sampling, which matters on real sensor buses even when accuracy gains are not visible in simulation.}
